\documentclass[11pt,a4paper]{article}
\usepackage[margin=22mm]{geometry}
\usepackage{fontspec}\setmainfont{Times New Roman}
\usepackage{amsmath,amssymb,amsthm,booktabs,array,rotating}
\usepackage[table]{xcolor}\definecolor{Bisque3}{RGB}{205,183,158}
\usepackage[hidelinks]{hyperref}
\usepackage{fancyhdr}\setlength{\headheight}{14pt}
\newtheorem{theorem}{Theorem}
\pagestyle{fancy}\fancyhf{}\fancyhead[L]{\small QEF: all-model bounds and source certification}
\fancyhead[R]{\small 19 October 2026}\fancyfoot[C]{\thepage}
\setlength{\parindent}{0pt}\setlength{\parskip}{5pt}
\let\subsubsection\section
\begin{document}
\begin{center}
{\LARGE\bfseries QEF: corrected Table 7 and source conditions}\\[5pt]
{\large Original three-layer operators; all-model reference calculations}\\[4pt]
{\small Numerical source certification is not established by the original generator specification.}
\end{center}
\setcounter{table}{6}
\subsubsection{Security model and a rotation-moment bound}
\label{sec:qef_moment_security}
{\color{red}
The security quantity considered here is the number of independently
encrypted quantum preparations needed to distinguish two specified images.
This is a sample-complexity statement, rather than a gate-complexity estimate
for exhaustive key search. The proof uses the original three QEF layers.
It does not assume that an adversary reconstructs an arbitrary unitary or
recovers every component of the original key.

\paragraph{Keys, ciphertexts, and the attack experiment.}
Write $A$ for the substitution-angle material produced by the nD-QHCS and
$B=(K^{\rm perm},K^{\rm ent})$ for the material produced by the
HG$\Phi$-QRNG. These are two source classes. The two components of $B$
are described by their joint distribution; independence is not inferred
from their output lengths. Each encryption uses a fresh secret key record,
and the receiver obtains the same record over an assumed secure channel.
Independent refresh between preparations is an explicit security
assumption, stronger than merely obtaining different numerical key values.
The experiment gives the adversary no additional side information
correlated with the target secret record; any claimed extension to
source leakage must bound that leakage explicitly.

Let $N=HW$, let $r$ denote the feature-register size, and let
$D=N2^r$ denote the dimension of the valid image subspace. Unused padded
position addresses are excluded from the ciphertext experiment. The image
states have the form
\begin{equation}
\label{eq:qef_security_image}
 |I\rangle=N^{-1/2}\sum_{k=0}^{N-1}|k\rangle|\chi_k\rangle,
 \qquad \langle\chi_k|\chi_k\rangle=1.
\end{equation}
The adversary is given the classical descriptions of two such images of
the same public dimensions. A uniform hidden bit $b$ selects $I_b$.
The adversary receives $q$ encryptions of that selected image, with
independently refreshed keys, and may perform any joint quantum
measurement. The publicly transmitted ciphertext is the quantum state,
not a classical list of all its amplitudes or a rendered image file.
Define
\begin{equation}
\label{eq:qef_averaged_cipher}
 \sigma_b=\mathbb E_{A,B}
   [U_{A,B}|I_b\rangle\langle I_b|U_{A,B}^{\dagger}].
\end{equation}
The joint ciphertext is $\sigma_b^{\otimes q}$. Its optimal binary
discrimination probability is \cite{Watrous2018QEF}
\begin{equation}
\label{eq:qef_operational_guess}
 P_{\rm guess}^{(q)}
 =\frac12+\frac14\left\|
       \sigma_0^{\otimes q}-\sigma_1^{\otimes q}\right\|_1.
\end{equation}
This experiment does not give repeated access to one fixed-key unitary.
It also does not cover key reuse, arbitrary reference-entangled chosen
inputs, authentication, or publication of simulator amplitude records.

\paragraph{Use the actual circuit structure.}
For the intended pixel-dependent substitution, write
\begin{equation}
 R_A=\sum_k |k\rangle\langle k|\otimes R_{A,k},\qquad
 R_{A,k}=\bigotimes_{j=1}^r T_{k,j},\quad
 T_{k,j}=R_z(\theta_{k,j}^z)R_y(\theta_{k,j}^y)R_x(\theta_{k,j}^x).
\end{equation}
Position-controlled CNOT/CZ gates leave each position basis label
unchanged. Thus, for suitable feature unitaries $V_B(t)$,
\begin{equation}
 E_B=\sum_t |t\rangle\langle t|\otimes V_B(t),\qquad
 U_{A,B}=E_BP_BR_A
 =\sum_k |\pi_B(k)\rangle\langle k|
       \otimes V_B(\pi_B(k))R_{A,k}.
\end{equation}
The permutation and entanglement keys may have any joint distribution
in the theorem below. In particular, a uniform distribution over all
$N!$ permutations is not needed for the image experiment in
Eq.~\eqref{eq:qef_security_image}. The permutation must nevertheless be
implemented reversibly with its work registers correctly cleaned.

\paragraph{An explicit angle reference law.}
The following law is an analytical reference, not an experimentally
certified property of the nD-QHCS. For an integer $L\geq2$, let each
Euler angle be uniform on
\begin{equation}
 \Theta_L=\{2\pi s/L:s=0,\ldots,L-1\}.
\end{equation}
Conditioned on $B$, the $3r$ angles within a pixel are independent,
and the angle tuples belonging to any two different pixels are
independent. Pairwise independence of pixel tuples is sufficient;
independence of all $N$ tuples jointly is not required. The continuous
reference is the corresponding independent uniform law on $[0,2\pi)$.
These conditions must hold after conditioning on any public setup
information or side information used in the attack experiment.
The reference is a key-law assumption, not a consequence of deterministic
seed expansion. For $N\geq2$, two independent uniform angle tuples alone
require at least $6r\log_2L$ bits of conditional source entropy. The
continuous law is an ideal limiting model; a finite implementation must
use the finite-law expression and provide sufficient actual key material.

Set
\begin{equation}
\label{eq:qef_beta_definition}
 a_L=\frac{1}{L\sin(\pi/(2L))},\qquad
 \beta_L=a_L^{6r},\qquad
 \beta_{\infty}=\left(\frac{2}{\pi}\right)^{6r}.
\end{equation}
Uniform Euler coordinates are not Haar-distributed rotations. The
nonzero first moment is retained in the analysis below.

\begin{theorem}[Conditional image-discrimination sample bound]
\label{thm:qef_moment_sample_bound}
Under the stated angle reference law and independent key refresh, let
$N\geq2$, $F=2^r$, $c_N=1-1/N$, and
\begin{equation}
\label{eq:qef_information_budget}
 \begin{split}
 x_+&=1+F\beta_L(N-1),\qquad x_-=1-F\beta_L,\\
 C_{N,r,L}
 &=\frac{x_+\ln x_+ +(N-1)x_-\ln x_-}{NF}.
 \end{split}
\end{equation}
Here $x_->0$ for the stated finite and continuous reference laws.
For every image state in Eq.~\eqref{eq:qef_security_image}, its average
ciphertext $\sigma$ satisfies
\begin{align}
 \frac12\|\sigma-I_D/D\|_1&\leq\beta_Lc_N,
 \label{eq:qef_one_image_distance}\\
 \mathsf D(\sigma\|I_D/D)&\leq C_{N,r,L},
 \label{eq:qef_cipher_relative_entropy}
\end{align}
where $\mathsf D$ uses natural logarithms. For any two such images,
achieving $P_{\rm guess}^{(q)}\geq p>1/2$ requires
\begin{equation}
\boxed{\displaystyle
 q\ \geq\left\lceil\max\left\{
 \frac{\ln2-h(p)}{C_{N,r,L}},\;
 \frac{p-1/2}{\beta_L(1-1/N)}
 \right\}\right\rceil,}
 \qquad h(p)=-p\ln p-(1-p)\ln(1-p).
\label{eq:qef_final_sample_bound}
\end{equation}
The continuous reference uses $\beta_\infty$ in the same expressions.
For $N=1$, the feature twirl makes the reference ciphertext exactly
maximally mixed, and no finite number of fresh encryptions distinguishes
the candidate images.
\end{theorem}

\begin{proof}
For one rotation axis $a$, averaging its conjugation over $\Theta_L$
dephases the Bloch components transverse to that axis. Composing the
independent $x$, $y$, and $z$ averages therefore gives
\begin{equation}
 \mathbb E[T X T^{\dagger}]=\frac{I_2}{2}\operatorname{Tr}X.
\end{equation}
For the first moment, a finite trigonometric sum gives
\begin{equation}
 \mathbb E[R_a(\theta)]
 =c_LI_2-is_L\sigma_a,
 \quad c_L=1/L,\quad s_L=L^{-1}\cot(\pi/(2L)).
\end{equation}
Each matrix on the right is $a_L$ times a unitary, since
$c_L^2+s_L^2=a_L^2$. Hence
$M=\mathbb E[R_{A,k}]=\beta_L^{1/2}W$ for a feature unitary $W$,
and the feature conjugation average is the completely depolarizing
channel. In the continuous limit,
\begin{equation}
 \mathbb E[T]=(2/\pi)^3 I_2,
 \qquad M=(2/\pi)^{3r}I_{2^r}.
\end{equation}
For example, the first identity follows by multiplying
$\mathbb E[R_a]=-i(2/\pi)\sigma_a$ in the order $z,y,x$.

Write $\rho=|I\rangle\langle I|=\sum_{k,l}|k\rangle\langle l|\otimes X_{kl}$,
and let $\Delta$ remove its off-diagonal position blocks. Conditional
pairwise independence gives
\begin{equation}
 \mathbb E[R_{A,k}X_{kl}R_{A,l}^{\dagger}]
 =\begin{cases}
 I_{2^r}\operatorname{Tr}(X_{kk})/2^r,&k=l,\\
 MX_{kl}M^{\dagger},&k\ne l.
 \end{cases}
\end{equation}
For Eq.~\eqref{eq:qef_security_image},
$\operatorname{Tr}X_{kk}=1/N$. Consequently the diagonal contribution
is $\tau=I_D/D$, which is invariant under every $E_BP_B$. The remaining
term is
\begin{equation}
 \sigma-\tau=\beta_L\,\mathbb E_B
  [F_B(\rho-\Delta\rho)F_B^{\dagger}],
 \quad F_B=E_BP_B(I\otimes W).
\label{eq:qef_off_diagonal_remainder}
\end{equation}
Every $F_B$ is unitary, so both trace norm and Hilbert--Schmidt norm
are bounded by the corresponding norms before this average.

The pinching map is an average over $N$ diagonal position-phase
unitaries, one of which is the identity. The triangle inequality gives
$\|\rho-\Delta\rho\|_1\leq2(1-1/N)$. Combining this with
Eq.~\eqref{eq:qef_off_diagonal_remainder} proves
Eq.~\eqref{eq:qef_one_image_distance}.
Moreover, because $\rho$ is pure and its position probabilities are
uniform,
\begin{equation}
 \|\rho-\Delta\rho\|_2^2
 =\operatorname{Tr}\rho^2-\operatorname{Tr}(\Delta\rho)^2
 =1-1/N.
\end{equation}
Thus $\|\sigma-\tau\|_2^2\leq\beta_L^2c_N$. A sharper entropy bound
uses the exact spectrum before the final average. The $N$ vectors
$|k\rangle|\chi_k\rangle$ are orthonormal, and on their span
$\rho-\Delta\rho$ has eigenvalues $1-1/N$ once and $-1/N$ with
multiplicity $N-1$. It is zero on the orthogonal complement.
Consequently $\tau+\beta_L(\rho-\Delta\rho)$ has eigenvalues
\begin{equation}
 \lambda_+=x_+/D,\qquad
 \lambda_-=x_-/D\quad(N-1\text{ times}),\qquad
 \lambda_0=1/D\quad(N(F-1)\text{ times}).
\end{equation}
Every conditional output in
Eq.~\eqref{eq:qef_off_diagonal_remainder} is a unitary conjugate of
this state. Relative entropy is convex and $\tau$ is invariant, so
\begin{equation}
 \mathsf D(\sigma\|\tau)
 \leq\lambda_+\ln(D\lambda_+)
       +(N-1)\lambda_-\ln(D\lambda_-)
 =C_{N,r,L}.
\end{equation}
For comparison, Jensen's inequality gives the valid, weaker bound
\begin{equation}
 C_{N,r,L}\leq\ln(1+D\beta_L^2c_N).
\end{equation}
For a uniform hidden bit and $q$ independent encryptions, the Holevo
information obeys
\begin{equation}
 \chi_q\leq\frac12\sum_{b=0}^1
 \mathsf D(\sigma_b^{\otimes q}\|\tau^{\otimes q})
 \leq qC_{N,r,L}.
\end{equation}
Indeed, subtracting $\chi_q$ from the first expression gives the
nonnegative relative entropy of the average ciphertext to
$\tau^{\otimes q}$. Holevo's theorem and binary Fano's inequality
\cite{Watrous2018QEF} imply that any guess with success at least $p$
requires mutual information at least $\ln2-h(p)$. Therefore
$\ln2-h(p)\leq qC_{N,r,L}$. Independently, the trace-distance triangle
inequality and Eq.~\eqref{eq:qef_one_image_distance} give
$\frac12\|\sigma_0-\sigma_1\|_1\leq2\beta_Lc_N$.
Expanding the difference of the tensor powers as a telescoping sum,
and using $\|\sigma_b\|_1=1$, yields
\begin{equation}
 \frac12\|\sigma_0^{\otimes q}-\sigma_1^{\otimes q}\|_1
 \leq 2q\beta_Lc_N,\qquad
 P_{\rm guess}^{(q)}\leq\frac12+q\beta_Lc_N.
\end{equation}
Success at least $p$ thus requires both inequalities inside the maximum
in Eq.~\eqref{eq:qef_final_sample_bound}. The proof allows arbitrary
joint measurements and unlimited computational postprocessing.
\end{proof}

\paragraph{Declared register conventions and numerical meaning.}
Table~\ref{tab:qef_quantum_attack_resistance} evaluates
Eq.~\eqref{eq:qef_final_sample_bound} at $p=2/3$ for all listed models
and four images. These are newly declared mathematical benchmark
settings; they are not reconstructed experimental parameters.
With $n=\lceil\log_2N\rceil$, the feature counts used here are
\begin{equation}
\begin{array}{c|c}
\text{Model}&r\\\hline
\mathrm{FRQI}&1\\
\mathrm{NEQR}&8\\
\mathrm{QIRHSI/EQIR}&10\\
\mathrm{AFQIRHSI}\ (p_{\rm AF}=4)&p_{\rm AF}+3=7\\
\mathrm{QIRBP}\ (b_{\rm BP}=0)&b_{\rm BP}+6=6\\
\mathrm{QRMW}\ (3\text{ channels})&8+\lceil\log_2 3\rceil=10\\
\mathrm{QIRWC}\ (b_{\rm WC}=0)&n+2b_{\rm WC}+8=n+8.
\end{array}
\end{equation}
All selected feature qubits are included in the controlled substitution
assumed by the theorem. In particular, QIRWC uses one main position
address of size $N$; its other position register is counted among the
feature qubits. Thus its total register count is $Q=2n+8$ and its valid
main-address dimension is $D=N2^{n+8}$. It is not counted as a
pair-address image with $N^2$ main addresses.

The stated QIRBP sum over positions, bit planes, and wavelength labels
has normalization
\[
 (N2^{b_{\rm BP}+3})^{-1/2}.
\]
The choice $b_{\rm BP}=0$ therefore
uses $r=6$ in this table. The previous illustrative $r=11$ calculation
would correspond to $b_{\rm BP}=5$ under this register convention; it
does not establish that this was an experimental setting.

The table values are necessary ciphertext-preparation counts under the
specified angle law and register conventions. They are neither measured
attack times nor bit-security exponents. The logarithm of a sample count
must not be labelled as elementary-gate security. Each displayed bound
is rounded conservatively; the accompanying calculation data retain
integer thresholds.

\begin{sidewaystable}[p]
\centering\color{red}
\caption{{\color{red}Conditional QEF image-discrimination sample bounds for all listed QIR models and the four original images. The original generator is not certified by these reference calculations.}}
\label{tab:qef_quantum_attack_resistance}
\vspace{0.25cm}
\textbf{(a) Ashurbanipal image, $H\times W=256\times256$, $n=\lceil\log_2(HW)\rceil=16$}\\[3pt]
{\footnotesize\setlength{\tabcolsep}{4pt}
\begin{tabular}{lccccccc}
\hline\rowcolor{Bisque3}
\textbf{Model} & $\boldsymbol{r}$ & $\boldsymbol{Q=n+r}$ & $\boldsymbol{D=N2^r}$ & $\boldsymbol{q_2}$ & $\boldsymbol{q_\infty}$ & $\boldsymbol{H_2\geq}$ & \textbf{Source}\\
\hline
\textbf{FRQI} & \(1\) & \(17\) & \(65{,}536\,2^{1}\) & $\geq 2$ & $\geq 3$ & \(6\) & NC \\
\textbf{NEQR} & \(8\) & \(24\) & \(65{,}536\,2^{8}\) & $\geq 2.79624\times10^{6}$ & $\geq 4.54809\times10^{10}$ & \(48\) & NC \\
\textbf{QIRHSI} & \(10\) & \(26\) & \(65{,}536\,2^{10}\) & $\geq 1.98605\times10^{9}$ & $\geq 5.77734\times10^{14}$ & \(60\) & NC \\
\textbf{EQIRHSI} & \(10\) & \(26\) & \(65{,}536\,2^{10}\) & $\geq 1.98605\times10^{9}$ & $\geq 5.77734\times10^{14}$ & \(60\) & NC \\
\textbf{AFQIRHSI} & \(7\) & \(23\) & \(65{,}536\,2^{7}\) & $\geq 349{,}531$ & $\geq 4.08705\times10^{8}$ & \(42\) & NC \\
\textbf{QIRBP} & \(6\) & \(22\) & \(65{,}536\,2^{6}\) & $\geq 43{,}692$ & $\geq 3.97626\times10^{6}$ & \(36\) & NC \\
\textbf{QRMW} & \(10\) & \(26\) & \(65{,}536\,2^{10}\) & $\geq 1.98605\times10^{9}$ & $\geq 5.77734\times10^{14}$ & \(60\) & NC \\
\textbf{QIRWC} & \(16+8\) & \(40\) & \(65{,}536\,2^{24}\) & $\geq 2.29734\times10^{30}$ & $\geq 3.12918\times10^{43}$ & \(144\) & NC \\
\hline\end{tabular}}\par

\par\vspace{0.3cm}
\textbf{(b) Hammurabi code image, $H\times W=222\times227$, $n=\lceil\log_2(HW)\rceil=16$}\\[3pt]
{\footnotesize\setlength{\tabcolsep}{4pt}
\begin{tabular}{lccccccc}
\hline\rowcolor{Bisque3}
\textbf{Model} & $\boldsymbol{r}$ & $\boldsymbol{Q=n+r}$ & $\boldsymbol{D=N2^r}$ & $\boldsymbol{q_2}$ & $\boldsymbol{q_\infty}$ & $\boldsymbol{H_2\geq}$ & \textbf{Source}\\
\hline
\textbf{FRQI} & \(1\) & \(17\) & \(50{,}394\,2^{1}\) & $\geq 2$ & $\geq 3$ & \(6\) & NC \\
\textbf{NEQR} & \(8\) & \(24\) & \(50{,}394\,2^{8}\) & $\geq 3.04381\times10^{6}$ & $\geq 5.91175\times10^{10}$ & \(48\) & NC \\
\textbf{QIRHSI} & \(10\) & \(26\) & \(50{,}394\,2^{10}\) & $\geq 2.57085\times10^{9}$ & $\geq 7.51324\times10^{14}$ & \(60\) & NC \\
\textbf{EQIRHSI} & \(10\) & \(26\) & \(50{,}394\,2^{10}\) & $\geq 2.57085\times10^{9}$ & $\geq 7.51324\times10^{14}$ & \(60\) & NC \\
\textbf{AFQIRHSI} & \(7\) & \(23\) & \(50{,}394\,2^{7}\) & $\geq 349{,}533$ & $\geq 5.29581\times10^{8}$ & \(42\) & NC \\
\textbf{QIRBP} & \(6\) & \(22\) & \(50{,}394\,2^{6}\) & $\geq 43{,}692$ & $\geq 5.05218\times10^{6}$ & \(36\) & NC \\
\textbf{QRMW} & \(10\) & \(26\) & \(50{,}394\,2^{10}\) & $\geq 2.57085\times10^{9}$ & $\geq 7.51324\times10^{14}$ & \(60\) & NC \\
\textbf{QIRWC} & \(16+8\) & \(40\) & \(50{,}394\,2^{24}\) & $\geq 2.98764\times10^{30}$ & $\geq 4.06943\times10^{43}$ & \(144\) & NC \\
\hline\end{tabular}}\par

\par\vspace{0.3cm}
\textbf{(c) Lamassu image, $H\times W=256\times256$, $n=\lceil\log_2(HW)\rceil=16$}\\[3pt]
{\footnotesize\setlength{\tabcolsep}{4pt}
\begin{tabular}{lccccccc}
\hline\rowcolor{Bisque3}
\textbf{Model} & $\boldsymbol{r}$ & $\boldsymbol{Q=n+r}$ & $\boldsymbol{D=N2^r}$ & $\boldsymbol{q_2}$ & $\boldsymbol{q_\infty}$ & $\boldsymbol{H_2\geq}$ & \textbf{Source}\\
\hline
\textbf{FRQI} & \(1\) & \(17\) & \(65{,}536\,2^{1}\) & $\geq 2$ & $\geq 3$ & \(6\) & NC \\
\textbf{NEQR} & \(8\) & \(24\) & \(65{,}536\,2^{8}\) & $\geq 2.79624\times10^{6}$ & $\geq 4.54809\times10^{10}$ & \(48\) & NC \\
\textbf{QIRHSI} & \(10\) & \(26\) & \(65{,}536\,2^{10}\) & $\geq 1.98605\times10^{9}$ & $\geq 5.77734\times10^{14}$ & \(60\) & NC \\
\textbf{EQIRHSI} & \(10\) & \(26\) & \(65{,}536\,2^{10}\) & $\geq 1.98605\times10^{9}$ & $\geq 5.77734\times10^{14}$ & \(60\) & NC \\
\textbf{AFQIRHSI} & \(7\) & \(23\) & \(65{,}536\,2^{7}\) & $\geq 349{,}531$ & $\geq 4.08705\times10^{8}$ & \(42\) & NC \\
\textbf{QIRBP} & \(6\) & \(22\) & \(65{,}536\,2^{6}\) & $\geq 43{,}692$ & $\geq 3.97626\times10^{6}$ & \(36\) & NC \\
\textbf{QRMW} & \(10\) & \(26\) & \(65{,}536\,2^{10}\) & $\geq 1.98605\times10^{9}$ & $\geq 5.77734\times10^{14}$ & \(60\) & NC \\
\textbf{QIRWC} & \(16+8\) & \(40\) & \(65{,}536\,2^{24}\) & $\geq 2.29734\times10^{30}$ & $\geq 3.12918\times10^{43}$ & \(144\) & NC \\
\hline\end{tabular}}\par

\par\vspace{0.3cm}
\textbf{(d) Babylon Lion image, $H\times W=768\times1024$, $n=\lceil\log_2(HW)\rceil=20$}\\[3pt]
{\footnotesize\setlength{\tabcolsep}{4pt}
\begin{tabular}{lccccccc}
\hline\rowcolor{Bisque3}
\textbf{Model} & $\boldsymbol{r}$ & $\boldsymbol{Q=n+r}$ & $\boldsymbol{D=N2^r}$ & $\boldsymbol{q_2}$ & $\boldsymbol{q_\infty}$ & $\boldsymbol{H_2\geq}$ & \textbf{Source}\\
\hline
\textbf{FRQI} & \(1\) & \(21\) & \(786{,}432\,2^{1}\) & $\geq 2$ & $\geq 3$ & \(6\) & NC \\
\textbf{NEQR} & \(8\) & \(28\) & \(786{,}432\,2^{8}\) & $\geq 2.79620\times10^{6}$ & $\geq 3.87854\times10^{9}$ & \(48\) & NC \\
\textbf{QIRHSI} & \(10\) & \(30\) & \(786{,}432\,2^{10}\) & $\geq 1.98872\times10^{8}$ & $\geq 4.81641\times10^{13}$ & \(60\) & NC \\
\textbf{EQIRHSI} & \(10\) & \(30\) & \(786{,}432\,2^{10}\) & $\geq 1.98872\times10^{8}$ & $\geq 4.81641\times10^{13}$ & \(60\) & NC \\
\textbf{AFQIRHSI} & \(7\) & \(27\) & \(786{,}432\,2^{7}\) & $\geq 349{,}526$ & $\geq 3.95317\times10^{7}$ & \(42\) & NC \\
\textbf{QIRBP} & \(6\) & \(26\) & \(786{,}432\,2^{6}\) & $\geq 43{,}691$ & $\geq 1.91498\times10^{6}$ & \(36\) & NC \\
\textbf{QRMW} & \(10\) & \(30\) & \(786{,}432\,2^{10}\) & $\geq 1.98872\times10^{8}$ & $\geq 4.81641\times10^{13}$ & \(60\) & NC \\
\textbf{QIRWC} & \(20+8\) & \(48\) & \(786{,}432\,2^{28}\) & $\geq 2.00742\times10^{35}$ & $\geq 4.22549\times10^{50}$ & \(168\) & NC \\
\hline\end{tabular}}\par
\vspace{0.2cm}
\parbox{\textwidth}{\footnotesize\color{red}
$q_2$ and $q_\infty$ are necessary independent ciphertext preparations for success $2/3$, from Eq.~\eqref{eq:qef_final_sample_bound}; the latter is an ideal continuous reference, not a finite-key claim.
$H_2$ is the necessary conditional source entropy (bits) for the $L=2$ independent-tuple reference, not a security exponent. NC means the original source conditions are not certified.
The declared benchmark uses 8-bit intensity, AFQIRHSI $p=4$, RGB QRMW, and one wavelength label ($b=0$) for QIRBP/QIRWC. In QIRWC the second position register is included among the features, so $r=n+8$.
These are explicit mathematical settings, not reconstructed experiment metadata.
Valid dimension $D$ differs from padded $d=2^Q$ when $N<2^n$.
Large lower bounds are rounded down; exact integer ceilings are supplied separately.}
\end{sidewaystable}


To relate the reference calculation to an actual generator, suppose
the resulting average ciphertext is within trace distance $\kappa$
of the reference ciphertext, uniformly over the allowed input images.
For example, a bound $\kappa$ on the total variation distance of the
actual joint finite angle/control-key law from its finite reference
law suffices. Comparing a finite law directly to a continuous law in
total variation is inappropriate; the reference must use the same
finite alphabet. Then the preceding proof gives
\begin{equation}
\label{eq:qef_source_robust_bound}
\begin{split}
 C_{\kappa}&=\min\!\left\{\ln D,\;
 \ln\!\left(1+D[\beta_L\sqrt{1-1/N}+2\kappa]^2\right)\right\},\\
 q&\geq\left\lceil\max\!\left\{
 \frac{\ln2-h(p)}{C_\kappa},\;
 \frac{p-1/2}{\beta_L(1-1/N)+\kappa}
 \right\}\right\rceil.
\end{split}
\end{equation}
The trace-distance bound becomes
$\frac12\|\sigma-\tau\|_1\leq\beta_L(1-1/N)+\kappa$.
This follows from $\|\sigma-\sigma_{\rm ref}\|_2
\leq\|\sigma-\sigma_{\rm ref}\|_1\leq2\kappa$.
The source bound must include correlations between pixels, axes,
source parameters, both components of $B$, and any relevant side
information. Statistical image metrics or balanced one-bit marginals
do not certify this joint bound.

The mapping $\theta=\gamma\bmod(q,\pi)$ does not by itself establish
the reference law. Even if $q\bmod\pi$ were uniform, its angle range
would be $[0,\gamma\pi)$; the reference continuous range requires
$\gamma=2$, as well as the stated joint moment conditions. No measured
value of $\kappa$ or proof of those conditions for the original
generators is supplied here. Therefore the numerical rows are
conditional benchmarks and are not asserted as certified resistance
of the experimental generator.

\paragraph{Why the range cannot be omitted.}
There is an explicit contrasting calculation within the same circuit
family. Suppose the Euler angles are instead independent uniform
variables on $[0,\pi)$, independently of $B$. Averaging the Bloch
rotations gives, for a data qubit initially in $|0\rangle$ or
$|1\rangle$,
\begin{equation}
 \langle X\rangle_0=(2/\pi)^2,\qquad
 \langle X\rangle_1=-(2/\pi)^2.
\end{equation}
Consider all-black and all-white images with a feature data qubit
initially $|0\rangle$ and $|1\rangle$, respectively. At output position
$t=0$, no position control is active, so $V_B(0)=I$. That position
occurs with probability $1/N$, regardless of the permutation.
An adversary can request fresh encryptions until measuring $t=0$,
then measure $X$ on that feature qubit. This stopping-time procedure
uses $N$ preparations on average and succeeds with probability
\begin{equation}
 \frac{1+(2/\pi)^2}{2}\simeq0.702642>2/3.
\end{equation}
For a fixed budget of $q$ preparations, use the first such outcome
and guess uniformly if none is obtained. The resulting success is
\begin{equation}
 P_{\rm attack}(q)=\frac12+\frac{2}{\pi^2}
      [1-(1-1/N)^q].
\end{equation}
Therefore this contrasting law admits a fixed-budget attack with
success at least $2/3$ at
\begin{equation}
 q_{\rm attack}=\left\lceil
 \frac{\ln(1-\pi^2/12)}{\ln(1-1/N)}\right\rceil.
\end{equation}
For the four images in table order these counts are $113,285$,
$87,111$, $113,285$, and $1,359,425$. They are attack upper bounds
under the contrasting law, distinct from the theorem's lower bounds
under its reference law. This is not a claimed attack on an
unspecified nD-QHCS distribution. It shows why the source range and
joint moments cannot be omitted from a numerical security claim.

{\color{red}
\paragraph{Certification of the original angle source.}
The reference law above must be distinguished from a property of the
implemented generator. Write $Z=q\bmod\pi$. If $q$ has density $f_q$,
then, for $0\leq z<\pi$,
\begin{equation}
 f_Z(z)=\sum_{m\in\mathbb Z}f_q(z+m\pi),\qquad
 f_\theta(t)=\gamma^{-1}f_Z(t/\gamma),\quad 0\leq t<\gamma\pi.
 \label{eq:qef_wrapped_source_density}
\end{equation}
Indeed, the inverse image of a measurable set $A\subset[0,\pi)$ is
the disjoint union of $A+m\pi$; summing their probabilities and changing
variables proves the first identity. Scaling proves the second.
Thus $\gamma=2$ and $f_Z=1/\pi$ are sufficient for the one-pass
continuous marginal reference on $[0,2\pi)$, but the modulo operation
does not establish either uniformity or independence. The manuscript
does not specify an actual numerical $\gamma$, a probability law for the
initial source state, or a sampling map that proves the required
conditional pixel/axis moments. Chaotic sensitivity and a bounded angle
range do not provide that missing probability law.

\paragraph{A bound stated directly in terms of source moments.}
Full independence is sufficient for the preceding reference calculation,
but is not necessary for a security bound. Let $F=2^r$, $D=NF$, and
condition on $B=b$ in its essential support. For the actual substitution
operators define
\begin{align}
 \delta_2&=\mathop{\mathrm{ess\,sup}}_{b,k,\|u\|=1}
 \left\|\mathbb E_{A\mid b}
 [R_{A,k}|u\rangle\langle u|R_{A,k}^{\dagger}]-I_F/F\right\|_2,
 \label{eq:qef_actual_diagonal_moment}\\
 \eta_2&=\mathop{\mathrm{ess\,sup}}_{b,k\ne l,\|u\|=\|v\|=1}
 \left\|\mathbb E_{A\mid b}
 [R_{A,k}|u\rangle\langle v|R_{A,l}^{\dagger}]\right\|_2.
 \label{eq:qef_actual_cross_moment}
\end{align}
Here vectors have Euclidean norm one and $\|\cdot\|_2$ denotes the
Hilbert--Schmidt matrix norm. Public setup information is also conditioned
on. These quantities account for dependence between
the two source classes, between axes, and between pixel tuples.

\begin{theorem}[Source-moment certification criterion]
\label{thm:qef_actual_source_moments}
For image states in Eq.~\eqref{eq:qef_security_image} and the original
clean three-layer unitary, the actual average ciphertext satisfies
\begin{equation}
 \mathsf D(\sigma\|I_D/D)\leq C_{\rm actual},\qquad
 C_{\rm actual}=\ln\!\left[1+F\delta_2^2+(N-1)F\eta_2^2\right].
 \label{eq:qef_actual_source_budget}
\end{equation}
If whole key records are independent between preparations, success
$p>1/2$ in the binary image experiment requires
\begin{equation}
 \boxed{\displaystyle q\geq
 \left\lceil\frac{\ln2-h(p)}{C_{\rm actual}}\right\rceil.}
 \label{eq:qef_actual_source_sample_bound}
\end{equation}
When $C_{\rm actual}=0$, no finite number of these preparations
distinguishes the candidate images.
\end{theorem}
\begin{proof}
Before applying $E_bP_b$, subtract $\tau=I_D/D$ from the conditional
average. Its diagonal position blocks are $1/N$ times the matrices in
Eq.~\eqref{eq:qef_actual_diagonal_moment}; its off-diagonal blocks are
$1/N$ times those in Eq.~\eqref{eq:qef_actual_cross_moment}.
Distinct blocks are orthogonal in the Hilbert--Schmidt inner product.
Unitary conjugation preserves this norm and $\tau$; convexity of its
square therefore gives
\begin{equation}
 \|\sigma-\tau\|_2^2
 \leq\delta_2^2/N+(1-1/N)\eta_2^2.
\end{equation}
Using $\mathsf D(\sigma\|\tau)\leq
\ln(1+D\|\sigma-\tau\|_2^2)$ proves
Eq.~\eqref{eq:qef_actual_source_budget}. The same Holevo/Fano argument
used above proves Eq.~\eqref{eq:qef_actual_source_sample_bound}.
No independence within an individual key record is used in this proof.
\end{proof}

The reference law has $\delta_2=0$ and $\eta_2=\beta_L$.
For an uncertified source, only the general inequalities
$\delta_2^2\leq1-1/F$ and $\eta_2\leq1$ are available. They give
$C_{\rm actual}\leq\ln D$ and only the trivial necessary count
$q\geq1$ for $p=2/3$ and $D\geq2$. This is not a successful attack;
it means the source information supplied does not certify a larger
count. A nontrivial value must follow from certified upper bounds on
the actual conditional moments, rather than being assigned from the
reference table.

\paragraph{Finite keys and the stronger independence assumption.}
Two independent uniform pixel tuples over an $L$-point angle alphabet
contain $6r$ independent symbols. Consequently
\begin{equation}
 H(A\mid B)\geq H(\Theta_k,\Theta_l\mid B)=6r\log_2L,
 \qquad k\ne l.
 \label{eq:qef_source_entropy_requirement}
\end{equation}
If all angles are deterministic expansions of an $s$-bit substitution
seed, their pair has at most $2^s$ possible values. For its support $S$,
the actual law gives probability one to $S$, whereas the independent
uniform pair law gives at most $2^s/L^{6r}$. Hence their total variation
distance is at least
\begin{equation}
 \max\{0,1-2^s/L^{6r}\}.
 \label{eq:qef_source_support_obstruction}
\end{equation}
For QIRBP ($r=11$) and a deterministic 128-bit substitution seed, this
lower bound is $15/16$ at $L=4$, and $1-2^{-400}$ at $L=256$.
These results concern the strong independent-symbol reference; they do
not exclude a correlated source satisfying the weaker operator-moment
criterion in Theorem~\ref{thm:qef_actual_source_moments}.

\paragraph{What the HG$\Phi$ measurement establishes.}
For the stated raw-bit model, an observer who knows the phase and
reference setting has single-bit guessing probability and conditional
pointwise min-entropy
\begin{equation}
 p_{\rm guess}(b\mid\phi,\delta)
 =\frac{1+|\cos(\phi+\delta)|}{2},\qquad
 H_{\min}(b\mid\phi,\delta)
 =-\log_2p_{\rm guess}(b\mid\phi,\delta).
 \label{eq:qef_raw_bit_source_entropy}
\end{equation}
The latter can be zero. A fixed deterministic coarse-graining or hash
cannot create missing entropy; a reversible map preserves it. Specifying
a randomness extractor would additionally require a source entropy
bound, seed assumptions, output length, and extraction error. None is
supplied here as a certification of the original generator.

\paragraph{Freshness alone does not supply the missing bound.}
An explicit countermodel uses two constant FRQI images with feature
states $|0\rangle$ and $|1\rangle$. Set the $x$ and $y$ rotation angles
to zero, refresh the $z$ angles independently for every encryption,
allow any position permutation, and let the entanglement component
select CZ gates. Every layer preserves the feature's $Z$ basis value.
One feature measurement therefore distinguishes the two images with
probability one using a single ciphertext. The whole key record is
fresh even though its gate-selection component is fixed. This is an
allowed source-law countermodel to an architecture-only claim, not an
attack attributed to the unspecified actual nD-QHCS/HG$\Phi$ law.

Finally, the supplied experimental implementation resets NumPy's
scalar-seeded pseudorandom generator and uses a fixed seed in its
driver. Its QIRBP branch performs shared channel masks and XOR chaining;
it does not implement the controlled complex Euler rotations or the
specified nD-QHCS/HG$\Phi$ generator. It therefore cannot certify the
moments in Eqs.~\eqref{eq:qef_actual_diagonal_moment}--
\eqref{eq:qef_actual_cross_moment}. The actual-source numerical
resistance remains unestablished; the preceding table entries are
conditional mathematical benchmarks.
\par}


\paragraph{Interpretation of the original attack labels.}
Independent known/chosen-image pairs from other operations do not
share the target key under the independent-refresh assumption. They
can still reveal a biased generator law or shared source parameters;
the conditional key law must therefore be specified. Generic key
search, if a reversible candidate verifier is available, has the
usual conditional $\Theta(\sqrt{M/t})$ query baseline for $M$ candidates
and $t$ marked candidates \cite{GDO,Boyer1998QEF,Zalka1999QEF}.
No particular old key-length value is used to certify QEF, and no
reduction proving that all attacks require exhaustive search is
claimed. The dimension or parameter count of a general quantum state
or unitary is not substituted for an attack lower bound.

Theorem~\ref{thm:qef_moment_sample_bound} concerns the mathematical
quantum operators and the stated key law. Classical array simulations
and their rendered images do not verify these quantum moment
conditions. In particular, the original scalar-seeded XOR/permutation
simulation is not an implementation of the controlled complex
Euler rotations assumed here. Image-quality and decryption metrics
are therefore reported as diagnostics, rather than as evidence that
the conditional security bound holds. The new probabilistic bound
is proved analytically; it is not included in the existing Lean~4
unitarity checks. Entanglement alone supplies neither authentication
nor a proof of resistance to active modification.
\par}
\clearpage
\section{Verification and effects of the corrections}
The numerical supplement has 128 rows: eight allocations, four native
image sizes, and four angle reference laws ($L=2,4,256$ and the ideal
continuous limit). Integer thresholds agree at 120 and 170 decimal
digits. The exact entropy denominator and the trace-distance converse
strengthen the earlier purity-based bound; they do not change encryption.

The main manuscript corrects both PSB phase signs and proves gauge
invariance. This affects the phase, Bernoulli probability, and any
keys generated from that probability; old source records need
recomputation. The original source classes are retained.
The security table keeps four panels and tan headers, separates
QRMW and QIRWC, and declares each register allocation. QIRBP's
plane-index qubits are counted separately from intensity precision.
The QIRWC secondary address is included in its feature register.

The experimental scope, statistical interpretations, source indexing,
resource counts, ancillary cleanup discussion, and numerical discussion
are adjusted consistently. Existing image metrics remain diagnostics
of the supplied array simulations. The source moment bound is an
analytic result, not part of the earlier Lean unitarity checks.
The original manuscript and the previous review copy are retained.

\paragraph{Independent verification of the strengthened bound.}
Four finite reference models were checked by exhaustive substitution-key enumeration (64, 512, 4096, and 4096 records), including complete superoperators, first moments, and feature twirls. Exact spectra and relative entropies were checked on 48 normalized image states. Maximum superoperator, spectrum, and entropy errors were 1.8e-16, 2.5e-16, and 5.83e-16, respectively. Optimal joint binary measurements in 10 reference cases satisfied both converses. Two eight-record sources with correlated substitution and control material were checked on 40 image states (block-orthogonality error 1.11e-16), with 5 additional joint-measurement checks. Eight source-perturbation cases checked the robustness inequalities. All 128 table rows, dimensions, and integer ceilings were independently recomputed at 145 decimal digits using a different pi algorithm and a cancellation-resistant entropy formula. The fixed-budget range attack and 48 fresh-record Z/CZ outcomes were also checked. These calculations support the derivation; they do not certify the original source law or report cryptanalytic attacks or hardware performance.

\clearpage
\begin{thebibliography}{99}
\bibitem{Watrous2018QEF} J. Watrous, The Theory of Quantum Information,
Cambridge University Press (2018), Chapters 3 and 5.
\url{https://cs.uwaterloo.ca/~watrous/TQI/}.
\bibitem{GDO} L. K. Grover, A fast quantum mechanical algorithm for database search, STOC (1996).
\bibitem{Boyer1998QEF} M. Boyer et al., Tight bounds on quantum searching,
Fortschritte der Physik 46 (1998). \url{https://arxiv.org/abs/quant-ph/9605034}.
\bibitem{Zalka1999QEF} C. Zalka, Physical Review A 60, 2746--2751 (1999).
\url{https://arxiv.org/abs/quant-ph/9711070}.
\bibitem{FRQI} P. Q. Le et al., Flexible representation of quantum images,
Quantum Information Processing (2011).
\bibitem{NEQR} Y. Zhang et al., NEQR, Quantum Information Processing (2013).
\bibitem{QIRHSI} G.-L. Chen et al., QIRHSI, Quantum Information Processing 21, 5 (2022).
\bibitem{EQIRHSI} M. Li et al., EQIRHSI, Quantum Information Processing 22, 334 (2023).
\bibitem{AFQIRHSI} N. A. Alwan et al., Scientific Reports 15, 40286 (2025).
\url{https://doi.org/10.1038/s41598-025-24168-4}.
\bibitem{QIRBPQIRWC} N. A. Alwan et al., AIMS Mathematics 10, 10994--11035 (2025).
\url{https://doi.org/10.3934/math.2025499}.
\bibitem{Rabei1999QEF} E. M. Rabei, Arvind, N. Mukunda, and R. Simon,
Bargmann invariants and geometric phases: a generalised connection,
Physical Review A 60, 3397--3409 (1999), Eq.~(2.4).
\url{https://arxiv.org/abs/quant-ph/9907080}.
\bibitem{QRMW2018QEF} E. \c{S}ahin and \.{I}. Y{\i}lmaz,
QRMW: Quantum representation of multi wavelength images,
Turkish Journal of Electrical Engineering and Computer Sciences 26,
768--779 (2018). \url{https://arxiv.org/abs/1801.10349}.
\end{thebibliography}
\end{document}